\documentclass[10pt,a4paper]{amsart}
\usepackage[margin=19mm]{geometry}
\usepackage{amsmath,amssymb,mathtools}
\usepackage[scr=rsfs,cal=euler]{mathalfa}
\usepackage{xfrac}
\usepackage[unicode,hidelinks,bookmarks=false]{hyperref}
\newcommand{\ie}{i.e.\ }
\newcommand{\cf}{cf.\ }
\newcommand{\resp}{respectively\ }
\newcommand{\Subsection}{\subsection}
\providecommand{\nobreakdash}{\nobreak\hbox{-}}
\setlength{\emergencystretch}{2em}
\multlinegap0pt
\usepackage{bm}
\usepackage{bbm}
\usepackage{dsfont}
\usepackage{xspace}
\usepackage{cancel}
\usepackage{mathtools}
\usepackage{amsthm}
\makeatletter
\@ifundefined{lemma}{\newtheorem{lemma}{Lemma}}{}
\makeatother
\newcommand{\E}{\mathbb{E}}
\newcommand{\N}{\mathbb{N}}
\renewcommand{\P}{\mathbb{P}}
\title[Strategies in POMDPs with Stage Duration]{Strategies in POMDPs with Stage Duration}
\author{Ivan Novikov}
\date{Source: arXiv:2603.16055v2 (CC BY 4.0)}
\begin{document}
\maketitle

\noindent\emph{Source and licence.} Strategies in POMDPs with Stage Duration. Version arXiv:2603.16055v2 (CC BY 4.0), distributed under the Creative Commons Attribution 4.0 International licence, as recorded in the arXiv licence metadata for this exact version and shown on the abstract page https://arxiv.org/abs/2603.16055v2. Full rights notice: licenses/arxiv-2603.16055-CC-BY-4.0.txt.\par\smallskip

\noindent\textit{Editorial scope.} Counted unit: the Lemma whose source label is \texttt{lemma1} in the article (source0.tex lines 373--377, source label \texttt{lemma1}), delivered together with the article's own complete original proof of it (source0.tex lines 523--751, one \texttt{proof} environment of 229 lines). No mathematics is written, repaired or re-derived: statement and proof are the article's own text line for line. The proof is detached from the statement in the source: the article states the result at lines 373--377 and prints its proof at lines 523--751, and the proof environment's own header names the statement, so the association is the article's own. Required context retained uncounted: aux-lemma (lines 492--508). Every definition, display and label of those lines is retained unchanged. Omitted from this package and not redistributed: 1--372, 378--491, 509--522, 752--1101. Nothing inside a retained excerpt is omitted or altered: every retained body is delivered exactly as the article's lines are written, so no omission receipt of any kind accompanies this package. Citations: the retained excerpts carry no citation command at all, so no bibliography entry of the article is transcribed and no reference is redistributed. Reference adaptation: none was needed and none was made; every reference inside the delivered unit resolves inside it, so no unresolved reference or citation is produced. Printed slips kept literal: no printed word, symbol, hypothesis or number of the counted statement or of its proof was altered. Layout and class: the article's own class is \texttt{article} with packages including cmap, babel, graphicx, floatflt, amsfonts, amsthm, amsmath, mathtools, gensymb, amssymb, color, wrapfig, natbib, url; the wrapper is the lane's pinned \texttt{amsart} editorial wrapper at 10pt on A4, which restates the theorem-like environments the retained text needs and adds the article's own macro definitions that the retained text needs (\texttt{\textbackslash E}, \texttt{\textbackslash N}, \texttt{\textbackslash P}), with extra wrapper packages (bm, bbm, dsfont, xspace, cancel, mathtools). The delivered numbering is the unit's own. Assets: no retained span contains a figure environment, a graphics-inclusion command, a PDF-page inclusion, a table or a TikZ picture; no image file is copied into this package, the package declares no asset dependency and no image file is redistributed with it. Licence: version arXiv:2603.16055v2 of this article is distributed under the Creative Commons Attribution 4.0 International licence, as recorded in the arXiv licence metadata for this exact version and shown on the abstract page https://arxiv.org/abs/2603.16055v2. A full rights notice is delivered at licenses/arxiv-2603.16055-CC-BY-4.0.txt. Characters: every retained excerpt is pure ASCII, so the delivered engine, XeLaTeX with its default Latin Modern text font, reports no missing character.
\begin{lemma}
Let $k \in \N_0, \omega \in \Omega$, and let: 
\begin{itemize}
	\item $\eta_{k+1}$ be a history of length $k+1$ in $G_1$;
	\item $\eta'_{k+1}$ be a history of length $T_{k+1}$ in $G_h$;
	\item $\mathcal H^{fil}_{k+1}$ be the filtered history of length $k+1$ (as defined previously);
	\item $\mathcal H^{rem}_{k+1}$ be the supplementary information in $\eta'_{k+1}$ not present in $\mathcal{H}_{k+1}^{fil}$.
\end{itemize}

Then $\mathcal{H}_{k+1}^{rem}$ is conditionally independent of $\omega'_{T_{k}+1}$ given $\mathcal H^{fil}_{k+1} = \eta_{k+1}$. That is,
	\begin{align*}
		\P_{\sigma}^h (\mathcal H^{rem}_{k+1} \mid \mathcal H^{fil}_{k+1} = \eta_{k+1}, \omega'_{T_k + 1} = \omega) &= \P_{\sigma}^h (\mathcal H^{rem}_{k+1} \mid \mathcal H^{fil}_{k+1} = \eta_{k+1})\\
		&\text{and}\\
		\P_{\sigma}^h(\omega'_{T_{k}+1} = \omega \mid \mathcal{H}_{k+1}^{rem}, \mathcal H^{fil}_{k+1} = \eta_{k+1}) &= \P_{\sigma}^h(\omega'_{T_{k}+1} = \omega \mid \mathcal H^{fil}_{k+1} = \eta_{k+1}).
	\end{align*}
\label{aux-lemma}
\end{lemma}

\begin{lemma}
For each $k \in \N^*$, we have
$$\E_{\widehat \sigma}^1 g(\omega_k, a_k) = \E_\sigma^h g(\omega_{T_k}, a_{T_k}).$$
\label{lemma1}
\end{lemma}

\begin{proof}[Proof of Lemma~\ref{lemma1}]
We will prove a more general statement: for any $\omega \in \Omega, a \in A, k \in \N^*$, and any history $\eta_k$ of length $k$, we have
\begin{equation}
\P_{\widehat \sigma}^1(\eta_k, \omega_k = \omega, a_k = a) = \P_{\sigma}^h(\mathcal H^{fil}_k = \eta_k, \omega'_{T_k} = \omega, a'_{T_k} = a).
\label{rooss10}
\end{equation}

This identity implies the lemma. Indeed, by \eqref{rooss10} we have
\begin{equation}
\begin{aligned}
 \P_{\widehat \sigma}^1(\omega_k = \omega, a_k = a) &= \sum_{\substack{\text{Histories } \eta_{k} \text{ of} \\ \text{length } k \text{ in } G_1}} \P_{\widehat \sigma}^1(\eta_k, \omega_k = \omega, a_k = a)
= \sum_{\substack{\text{Histories } \eta_{k} \text{ of} \\ \text{length } k \text{ in } G_1}} \P_{\sigma}^h(\mathcal H^{fil}_k = \eta_k, \omega'_{T_k} = \omega, a'_{T_k} = a) \\
&= \P_{\sigma}^h(\omega'_{T_k} = \omega, a'_{T_k} = a).
\end{aligned}
\label{rooss101}
\end{equation}

Consequently, we obtain by \eqref{rooss101}
$$\E_{\widehat \sigma}^1 g(\omega_k, a_k) =
\sum_{\omega \in \Omega, a \in A}\P_{\widehat \sigma}^1(\omega_k = \omega, a_k = a) \cdot g(\omega, a) 
= \sum_{\omega \in \Omega, a \in A} \P_{\sigma}^h(\omega'_{T_k} = \omega, a'_{T_k} = a) \cdot g(\omega, a)
= \E_\sigma^h g(\omega'_{T_k}, a'_{T_k}).$$


The rest of the proof is dedicated to justifying \eqref{rooss10}. The proof proceeds by induction on $k$. We first consider the base case $k=1$. Since the initial probability distribution is the same in both $G_1$ and $G_h$, we have
\begin{equation}
\P_{\widehat \sigma}^1(\eta_1, \omega_1 = \omega) = \P_{\sigma}^h(\mathcal H^{fil}_1 = \eta_1, \omega_1 = \omega).
\label{tralala1}
\end{equation}

Because the state is frozen between stage $T_k + 1$ and stage $T_{k+1}$, we have for any $k$
\begin{equation}
\omega'_{T_{k+1}} = \omega'_{T_{k}+1}.
\label{roo10x3}
\end{equation}


By the definition of conditional probability, we have
\begin{equation}
\P_{\widehat \sigma}^1(\eta_1, \omega_1 = \omega, a_1 = a) = \P_{\widehat \sigma}^1(a_1 = a \mid \eta_1, \omega_1 = \omega) \cdot \P_{\widehat \sigma}^1(\eta_1, \omega_1 = \omega) = \widehat \sigma(\eta_1)(a) \cdot \P_{\widehat \sigma}^1(\eta_1, \omega_1 = \omega).
\label{tralala2}
\end{equation}

By the definition of conditional probability and by \eqref{roo10x3}, we have
\begin{equation}
\begin{aligned}
\P_{\sigma}^h(\mathcal H^{fil}_1 = \eta_1, \omega'_{T_1} = \omega, a'_{T_1} = a) 
&= \P_{\sigma}^h(a'_{T_1} = a \mid \mathcal H^{fil}_1 = \eta_1, \omega'_{T_1} = \omega) \cdot \P_{\sigma}^h(\mathcal H^{fil}_1 = \eta_1, \omega'_{T_1} = \omega) \\
&= \P_{\sigma}^h(a'_{T_1} = a \mid \mathcal H^{fil}_1 = \eta_1, \omega'_{1} = \omega) \cdot \P_{\sigma}^h(\mathcal H^{fil}_1 = \eta_1, \omega'_{1} = \omega).
\end{aligned}
\label{tralala3}
\end{equation}

Now, let us fix a history $\eta'_{1}$ of length $T_1$ in $G_h$. Provided that the event 
$$\left\{\mathcal H^{fil}_1 = \eta_1, \eta'_{1}, \omega'_{1} = \omega\right\}$$
has positive probability, we have $\eta'_{1} = \left(\mathcal H^{fil}_1, \mathcal H^{rem}_1\right)$. Thus, by the law of total probability, we obtain
\begin{equation}
\begin{aligned}
&\P_{\sigma}^h (a'_{T_1} = a \mid \mathcal H^{fil}_1 = \eta_1, \omega'_{1} = \omega) 
\\
=&
\sum_{\substack{\text{Histories } \eta'_{1} \text{ of} \\ \text{length } T_{1} \text{ in } G_h}}
\P_{\sigma}^h (a'_{T_1} = a \mid \mathcal H^{fil}_1 = \eta_1, \eta'_{1}, \omega'_{1} = \omega) \cdot 
\P_{\sigma}^h (\eta'_{1} \mid \mathcal H^{fil}_1 = \eta_1, \omega'_{1} = \omega) \\
=&
\sum_{\mathcal H^{rem}_1}
\P_{\sigma}^h (a'_{T_1} = a \mid \mathcal H^{fil}_1 = \eta_1, \mathcal H^{rem}_1, \omega'_{1} = \omega) \cdot 
\P_{\sigma}^h (\mathcal H^{rem}_1 \mid \mathcal H^{fil}_1 = \eta_1, \omega'_{1} = \omega).
\end{aligned}
\label{tralala4}
\end{equation}

Since $a'_{T_{1}}$ depends by definition only on the history of length $T_{1}$, we have
\begin{equation}
\P_{\sigma}^h (a'_{T_1} = a \mid \mathcal H^{fil}_1 = \eta_1, \mathcal H^{rem}_1, \omega'_{1} = \omega) = \P_{\sigma}^h (a'_{T_1} = a \mid \mathcal H^{fil}_1 = \eta_1, \mathcal H^{rem}_1).
\label{tralala5}
\end{equation}


By Lemma~\ref{aux-lemma}, we have
\begin{equation}
\P_{\sigma}^h (\mathcal H^{rem}_1 \mid \mathcal H^{fil}_1 = \eta_1, \omega'_{1} = \omega) = \P_{\sigma}^h (\mathcal H^{rem}_1 \mid \mathcal H^{fil}_1 = \eta_1).
\label{tralala6}
\end{equation}

From \eqref{tralala4}, \eqref{tralala5}, \eqref{tralala6}, and the law of total probability, we obtain
\begin{equation}
\begin{aligned}
\P_{\sigma}^h (a'_{T_1} = a \mid \mathcal H^{fil}_1 = \eta_1, \omega'_{1} = \omega) 
&=
\sum_{\mathcal H^{rem}_1}
\P_{\sigma}^h (a'_{T_1} = a \mid \mathcal H^{fil}_1 = \eta_1, \mathcal H^{rem}_1) \cdot 
\P_{\sigma}^h (\mathcal H^{rem}_1 \mid \mathcal H^{fil}_1 = \eta_1) \\
&=
\P_{\sigma}^h (a'_{T_1} = a \mid \mathcal H^{fil}_1 = \eta_1)
=
\widehat \sigma(\eta_1)(a).
\end{aligned}
\label{tralala7}
\end{equation}

From \eqref{tralala1}, \eqref{tralala2}, \eqref{tralala3}, and \eqref{tralala7}, we obtain
$$\P_{\widehat \sigma}^1(\eta_1, \omega_1 = \omega, a_1 = a) = \P_{\sigma}^h(\mathcal H^{fil}_1 = \eta_1, \omega'_{T_1} = \omega, a'_{T_1} = a). $$
This completes the proof of the base case for the induction.

We now suppose that \eqref{rooss10} holds for some $k \in \N^*$, and we will prove that it also holds for $k+1$. Let $\eta_{k+1} = (\eta_k, \overline a, \overline s).$

In $G_1$, we have (similarly to \eqref{tralala2})
\begin{equation} 
\P_{\widehat \sigma}^1(\eta_{k+1}, \omega_{k+1} = \omega, a_{k+1} = a) = \P_{\widehat \sigma}^1(\eta_{k+1}, \omega_{k+1} = \omega) \cdot \widehat \sigma(\eta_{k+1})(a).
\label{roo10}
\end{equation}

By construction of the events $\mathcal H^{fil}_{k+1} = \eta_{k+1}$ and $\eta_{k+1}$, we can write them as unions of disjoint events:

\begin{equation}
\begin{aligned}
\left\{\eta_{k+1}\right\} &= \bigcup_{\widehat \omega \in \Omega} \left\{\eta_{k}, \omega'_{T_{k}} = \widehat \omega, a'_{T_{k}} = \overline a, s'_{T_{k}+1} = \overline s \right\}; \\
\left\{\mathcal H^{fil}_{k+1} = \eta_{k+1}\right\} &= \bigcup_{\widehat \omega \in \Omega} \left\{\mathcal H^{fil}_{k} = \eta_{k}, \omega'_{T_{k}} = \widehat \omega, a'_{T_{k}} = \overline a, s'_{T_{k}+1} = \overline s \right\}.
\end{aligned}
\label{roo10x4}
\end{equation}


By \eqref{roo10x4}, the law of total probability, and the definition of conditional probability, we have
\begin{equation}
\begin{aligned}
\mathbb{P}_{\widehat{\sigma}}^1(\eta_{k+1}, \omega_{k+1} = \omega) 
&= \sum_{\widehat{\omega} \in \Omega} \mathbb{P}_{\widehat{\sigma}}^1(\eta_k, \omega_k = \widehat{\omega}, a_k = \overline{a}, \omega_{k+1} = \omega, s_{k+1} = \overline{s}) \\
&= \sum_{\widehat{\omega} \in \Omega} \mathbb{P}_{\widehat{\sigma}}^1(\eta_k, \omega_k = \widehat{\omega}, a_k = \overline{a}) \cdot \mathbb{P}_{\widehat{\sigma}}^1(\omega_{k+1} = \omega, s_{k+1} = \overline{s} \mid \eta_k, \omega_k = \widehat{\omega}, a_k = \overline{a}) \\
&= \sum_{\widehat{\omega} \in \Omega} \mathbb{P}_{\widehat{\sigma}}^1(\eta_k, \omega_k = \widehat{\omega}, a_k = \overline{a}) \cdot P(\omega \mid \overline{a}, \widehat{\omega}) \cdot \mathbf{1}_{f(\omega) = \overline{s}}
\end{aligned}
\label{roo10x1}
\end{equation}

In $G_h$, it follows from the definition of conditional probability that
\begin{multline}
\P_{\sigma}^h(\mathcal H^{fil}_{k+1} = \eta_{k+1}, \omega'_{T_{k+1}} = \omega, a'_{T_{k+1}} = a) = \\ \P_{\sigma}^h(\mathcal H^{fil}_{k+1} = \eta_{k+1}, \omega'_{T_{k+1}} = \omega) \cdot \P_{\sigma}^h(a'_{T_{k+1}} = a \mid \mathcal H^{fil}_{k+1} = \eta_{k+1}, \omega'_{T_{k+1}} = \omega).
\label{roo11}
\end{multline}

Using \eqref{roo10x3}, \eqref{roo10x4}, the law of total probability, and the definition of conditional probability, we obtain
\begin{equation}
\begin{aligned}
&\P_{\sigma}^h(\mathcal H^{fil}_{k+1} = \eta_{k+1}, \omega'_{T_{k+1}} = \omega) = \sum_{\widehat \omega \in \Omega}
\P_{\sigma}^h(\mathcal H^{fil}_{k} = \eta_{k}, \omega'_{T_{k}+1} = \omega, \omega'_{T_{k}} = \widehat \omega, a'_{T_{k}} = \overline a, s'_{T_{k}+1} = \overline s) \\
&=
\sum_{\widehat \omega \in \Omega}
\P_{\sigma}^h(\mathcal H^{fil}_{k} = \eta_{k}, \omega'_{T_{k}} = \widehat \omega, a'_{T_{k}} = \overline a) \cdot 
\P_{\sigma}^h(\omega'_{T_{k}+1} = \omega, s'_{T_{k}+1} = \overline s \mid \mathcal H^{fil}_{k} = \eta_{k}, \omega'_{T_{k}} = \widehat \omega, a'_{T_{k}} = \overline a). \\
&=
\sum_{\widehat \omega \in \Omega}
\P_{\sigma}^h(\mathcal H^{fil}_{k} = \eta_{k}, \omega'_{T_{k}} = \widehat \omega, a'_{T_{k}} = \overline a) \cdot P(\omega \mid \overline a, \widehat \omega) \cdot \mathbf{1}_{f(\omega) = \overline s}
\end{aligned}
\label{roo10x2}
\end{equation}


Thus, by \eqref{roo10x1}, \eqref{roo10x2}, and the induction hypothesis, we have
\begin{equation}
\mathbb{P}_{\widehat{\sigma}}^1(\eta_{k+1}, \omega_{k+1} = \omega) = \P_{\sigma}^h(\mathcal H^{fil}_{k+1} = \eta_{k+1}, \omega'_{T_{k+1}} = \omega)
\label{roo08}
\end{equation}

By the law of total probability, we obtain
\begin{multline}
\P_{\sigma}^h(a'_{T_{k+1}} = a \mid \mathcal H^{fil}_{k+1} = \eta_{k+1}, \omega'_{T_{k}+1} = \omega) = \\ 
\sum_{\substack{\text{Histories } \eta'_{k+1} \text{ of} \\ \text{length } T_{k+1} \text{ in } G_h}}
\P_{\sigma}^h(a'_{T_{k+1}} = a \mid \eta'_{k+1}, \mathcal H^{fil}_{k+1} = \eta_{k+1}, \omega'_{T_{k}+1} = \omega) \cdot \P_{\sigma}^h(\eta'_{k+1} \mid \mathcal H^{fil}_{k+1} = \eta_{k+1}, \omega'_{T_{k}+1} = \omega),
\label{roo04}
\end{multline}

Since $a'_{T_{k+1}}$ depends by definition only on the history of length $T_{k+1}$, we have

\begin{equation}
\P_{\sigma}^h(a'_{T_{k+1}} = a \mid \eta'_{k+1}, \mathcal H^{fil}_{k+1} = \eta_{k+1}, \omega'_{T_{k}+1} = \omega) = \P_{\sigma}^h(a'_{T_{k+1}} = a \mid \eta'_{k+1}).
\label{roo03}
\end{equation}

By the Bayes' theorem, we obtain
\begin{equation}
\P_{\sigma}^h(\eta'_{k+1} \mid \mathcal H^{fil}_{k+1} = \eta_{k+1}, \omega'_{T_{k}+1} = \omega) = 
\frac{\P_{\sigma}^h(\omega'_{T_{k}+1} = \omega \mid \eta'_{k+1}, \mathcal H^{fil}_{k+1} = \eta_{k+1}) \cdot \P_{\sigma}^h(\eta'_{k+1} \mid \mathcal H^{fil}_{k+1} = \eta_{k+1})}{\P_{\sigma}^h(\omega'_{T_{k}+1} = \omega \mid \mathcal H^{fil}_{k+1} = \eta_{k+1})}.
\label{roo01}
\end{equation}

%Note that $\omega'_{T_{k}+1}$ is conditionaly independent with 

Provided that the event 
$$\{\eta'_{k+1}, \mathcal H^{fil}_{k+1} = \eta_{k+1}\}$$
 has positive probability, we have $\eta'_{k+1} = \left(\mathcal{H}_{k+1}^{fil}, \mathcal{H}_{k+1}^{rem}\right)$. Consequently, we obtain
\begin{equation}
\P_{\sigma}^h(\omega'_{T_{k}+1} = \omega \mid \eta'_{k+1}, \mathcal H^{fil}_{k+1} = \eta_{k+1}) =
\P_{\sigma}^h(\omega'_{T_{k}+1} = \omega \mid \mathcal{H}_{k+1}^{rem}, \mathcal H^{fil}_{k+1} = \eta_{k+1})
\label{rttt02}
\end{equation}

By Lemma~\ref{aux-lemma}, we have

\begin{equation}
\P_{\sigma}^h(\omega'_{T_{k}+1} = \omega \mid \mathcal{H}_{k+1}^{rem}, \mathcal H^{fil}_{k+1} = \eta_{k+1}) =
\P_{\sigma}^h(\omega'_{T_{k}+1} = \omega \mid \mathcal H^{fil}_{k+1} = \eta_{k+1}),
\label{rttt01}
\end{equation}

Thus, by \eqref{roo01}, \eqref{rttt02}, and \eqref{rttt01}, we obtain

\begin{equation}
\P_{\sigma}^h(\eta'_{k+1} \mid \mathcal H^{fil}_{k+1} = \eta_{k+1}, \omega'_{T_{k}+1} = \omega) = 
\P_{\sigma}^h(\eta'_{k+1} \mid \mathcal H^{fil}_{k+1} = \eta_{k+1}).
\label{roo02}
\end{equation}

From \eqref{roo04}, \eqref{roo03}, \eqref{roo02}, and the law of total probability, we have
\begin{equation}
\begin{aligned}
\P_{\sigma}^h(a'_{T_{k+1}} = a \mid \mathcal H^{fil}_{k+1} = \eta_{k+1}, \omega'_{T_{k}+1} = \omega) &=  \sum_{\substack{\text{Histories } \eta'_{k+1} \text{ of} \\ \text{length } T_{k+1} \text{ in } G_h}}
\P_{\sigma}^h(a'_{T_{k+1}} = a \mid \eta'_{k+1}) \cdot \P_{\sigma}^h(\eta'_{k+1} \mid \mathcal H^{fil}_{k+1} = \eta_{k+1}) \\
&= \P_{\sigma}^h(a'_{T_{k+1}} = a \mid  \mathcal H^{fil}_{k+1} = \eta_{k+1}) = \widehat \sigma(\eta_{k+1})(a).
\end{aligned}
\label{roo09}
\end{equation}

Finally, by combining \eqref{roo10}, \eqref{roo11}, \eqref{roo08}, and \eqref{roo09}, we have
$$\P_{\widehat \sigma}^1(\eta_{k+1}, \omega_{k+1} = \omega, a_{k+1} = a) = 
\P_{\sigma}^h(\mathcal H^{fil}_{k+1} = \eta_{k+1}, \omega'_{T_{k+1}} = \omega, a'_{T_{k+1}} = a).$$

This completes the proof of the induction step.
\end{proof}

\end{document}
